\printedpage{84}
% head: A. Grothendieck : Dérivateurs, Ch. VII, Édition des Universités de Montpellier II et Paris VII
% foot: Version 22 décembre 2001

i.e.
\[
\underbrace{\hot(\uHom(Y,C))\times\hot(Y)}_{\hot(\uHom(Y,C)\times Y)}\longrightarrow\hot(C),
\]
which arrow is simply deduced from
\[
\uHom(Y,C)\times Y\longrightarrow C.
\]
But the arrow $(*)$ has no desire to be \texttt{[an]} isomorphism. It is because it is already not one in the special case $Y=\mathbf{\Delta}^1$, essentially, which forces us, in order to make a theory of paths in categories that is topologically reasonable, not to confine ourselves to $\mathbf{\Delta}^1$ as a type of paths, but to take composite intervals $\tau\in\mathrm{Ty}$.

\msp{page 136}

On the other hand, $(*)$ is \texttt{[an]} isomorphism when $C$ is a Thomason-Kan category. It suffices to see that for every $Z$ in (Cat), the Hom from $\hot(Z)$ into the two sides is isomorphic via the arrow deduced from $(*)$. Now if $C$ is Thomason-Kan, so is $\uHom(Y,C)$ (immediate from the definition of Thomason-Kan in terms of extensions of
\begin{tikzcd}[baseline=(\tikzcdmatrixname-1-1.base), row sep=small]
A_0 \arrow[r, hook, "i"] \arrow[dr] & A \arrow[d, dotted] \\
& C
\end{tikzcd}
$\ldots$(\footnote[29]{a bit hasty! One has to know that the cofibrant maps $i$ of \textsc{Thomason} remain so after cartesian multiplication by $Y$ \ldots})), hence
\[
\mathrm{Hom}_{\mathrm{Hot}}(\hot(Z),\hot(\uHom(Y,C)))\xleftarrow{\ \sim\ }\pi_0(\underbrace{\uHom_{\mathrm{Cat}}(Z,\uHom(Y,C))}_{\uHom_{\mathrm{Cat}}(Z\times Y,C)})\simeq\pi_0\uHom(Z\times Y,C), % Reading resolved against the supplied PDF; see corrections.csv.
\]
and
\[
\begin{aligned}
&\mathrm{Hom}_{\mathrm{Hot}}(\hot(Z),\uHom_{\mathrm{Hot}}(\hot(Y),\hot(C)))\simeq\mathrm{Hom}_{\mathrm{Hot}}(\underbrace{\hot(Z)\times\hot(Y)}_{\hot(Z\times Y)},\hot(C))\\
&\qquad\underbrace{\xleftarrow{\ \sim\ }}_{\text{since $C$ is Thomason-Kan}}\qquad\quad\pi_0(\uHom_{\mathrm{Cat}}(Z\times Y,C)),
\end{aligned}
\]
O.K.

We thus see that the categories $\uCh_{\infty,\Sigma}(X;a,b)$ have no tendency to be Thomason-Kan categories. Yet, its type in Hot must be ``the right one'', without first having to replace $X$ by a Thomason-Kan envelope $\bar{X}$.

\printedpage{85}
% head: A. Grothendieck : Dérivateurs, Ch. VII, Édition des Universités de Montpellier II et Paris VII
% foot: Version 22 décembre 2001

\msp{page 137}

\section*{12\quad Back to the ``categories of paths'' in \textmd{Cat}, as a category of models for \textmd{Hot}}

I see exactly two interesting applications of the formalism of paths in categories. One, which above all has forced my hand, is to be able to construct Hom\,ext, \textit{external} Hom with values in Hot, for every prederivator $\mathbf{D}$ defined by a pair $(\mathcal{M},\Sigma)$ (and I know of no prederivator that is not defined this way).

The other application is to have in (Cat) a \textit{functorial construction,} in terms of an arbitrary object $X$ or an arrow $X\longrightarrow Y$, of the analog of the formalism of \textsc{Cartan-Serre} on path spaces. In this setting, it is understood that one always takes $\Sigma_X=\mathrm{Fl}(X)$, the set of \textit{all} arrows of $X$. To construct functorially in $X$ in (Cat)
\[
\begin{tikzcd}[row sep=huge]
\uCh_\infty(X) \arrow[d] \\
X\times X,
\end{tikzcd}
\]
one must have fixed an inverse system
\[
\underline{t}=(t_n)_{n\in\mathbf{N}}
\]
\msp{page 138}
in \underline{Ty} which is cofinal in \underline{Ty}, which simply means that among the $t_n$ there are some of arbitrarily large useful length (and since
\[
n\longmapsto\mathrm{long.ut.}(t_n)
\]
is increasing, this implies that long.ut.$(t_n)$ tends to $+\infty$ for $n\longrightarrow+\infty$). The main facts to check next (if it can be done) are the essential properties of the arrows in the diagram
\printedpage{86}
% head: A. Grothendieck : Dérivateurs, Ch. VII, Édition des Universités de Montpellier II et Paris VII
% foot: Version 22 décembre 2001
\[
\begin{tikzcd}[column sep=large, row sep=large]
& X \arrow[d, "\tau"] & \\
& \uCh_\infty(X) \arrow[d, "p"] \arrow[ddl, "\sigma"'] \arrow[ddr, "\beta"] & \\
& X\times X \arrow[dl, "\mathrm{pr}_1"] \arrow[dr, "\mathrm{pr}_2"'] & \\
X & & X\ ,
\end{tikzcd}
\qquad\qquad
\left|\begin{array}{l}
\sigma\tau=\beta\tau=\mathrm{id}_X\\
p\tau=\mathrm{diag}_X
\end{array}\right.
\]
and more generally, for an arrow
\[
X'\xrightarrow{\ f\ }X
\]
in Cat, of the arrows in the corresponding diagram
\[
\begin{tikzcd}
& \overline{X'}=C(f) \arrow[dl] \arrow[dr, "g"] \arrow[ddrr, bend left=30, "\bar{f}"] & & \\
X' \arrow[ur, dotted, bend left=40, "\alpha"] \arrow[dr, "f"'] & |[draw, rounded corners]| \text{cart.} & \uCh_\infty(X) \arrow[dl, "\sigma"] \arrow[dr, "\beta"] & \\
& X & & X,
\end{tikzcd}
\qquad
\bar{f}\alpha=(\beta g)\alpha=f,
\]
and more particularly of $\beta g$. The arrow
\msp{page 139}
in dotted lines $\alpha$ is the section of $\overline{X'}$ over $X'$ deduced from the section $\tau$ of $\uCh_\infty(X)$ over $X$ which realizes the decomposition of $f$ into
\[
X'\xrightarrow{\ \alpha\ }\overline{X'}\xrightarrow{\ \bar{f}\ }X,
\]
where $\bar{f}$ should be a \textit{``fibrant''} arrow and a Hot-equivalence, $\alpha$ a \textit{``cofibrant''} arrow in the sense of \textsc{Thomason}. (But whether this is so, remains to be seen!) Here for the moment are the essential properties I would like to check.

\begin{enumerate}[label={}, leftmargin=*]
\item[1$^o$)] $p$ is smooth and proper, and is a HOT-fibration.

\item[2$^o$)] So the same holds for $\sigma=\mathrm{pr}_1p$ and $\beta=\mathrm{pr}_2p$. Moreover, these HOT-fibrations are \textit{trivial,} i.e. with aspheric fibers.

\printedpage{87}
% head: A. Grothendieck : Dérivateurs, Ch. VII, Édition des Universités de Montpellier II et Paris VII
% foot: Version 22 décembre 2001

\item[3$^o$)] More generally (as regards $\beta$), $\overline{X'}\xrightarrow{\ \bar{f}\ }X$ should still be proper and smooth and a HOT-fibration, but of course in general not trivial. (It should be trivial if and only if $f$ is a Hot-equivalence.) (\footnote[30]{follows formally from 1$^o$ and 2$^o$, cf. page 170.}).
\end{enumerate}

A word on ``duality''. The statements about $\sigma$ (resp. $\beta$) translate into statements about $\beta$ (resp. $\sigma$) by passing from $\underline{t}$ to $\underline{t}^o=(t^o_n)$, where $\tau\longmapsto\tau^o$ is the functor in \underline{Ty} ``passage to the opposite interval''. We have in fact % Reading resolved against the supplied PDF; see corrections.csv.

\msp{page 140}

an isomorphism functorial in $X$ (passage to the opposite path)
\[
\boxed{\begin{array}{ccc}
c & \longmapsto & c^o\\[4pt]
\uCh^{\underline{t}}_\infty(X) & \longrightarrow & \uCh^{\underline{t}^o}_\infty(X)
\end{array}}
\qquad\qquad
\text{(\footnotemark[31])}
\]
\footnotetext[31]{cf. page 143 for the passage from $\underline{t}$ to $\underline{t}^-$.}
which ``exchanges source and target''. Of course, we do well to take $\underline{t}$ such that we have $\underline{t}=\underline{t}^o$ (this is what I had done, by a sound instinct, right from the start in V), so that we have a canonical involutive functor
\[
\begin{array}{ccc}
c & \longmapsto & c^o\\[4pt]
\uCh_\infty(X) & \longrightarrow & \uCh_\infty(X)
\end{array}
\]
(I again drop the superscript $\underline{t}$), which ``exchanges source and target''. So we find, up to canonical isomorphism, the ``same'' factorization of $f:X'\longrightarrow X$ into $\bar{f}\circ\alpha$, if we do as on page 38 to define $\overline{X'}$ and $\alpha$, or if in this diagram we interchange the roles of $\sigma$ and $\beta$. % [sic] page 38 -> page 138, the diagram defining \overline{X'} and \alpha stands on manuscript page 138 (apparently)

Finally:

\begin{enumerate}[label={}, leftmargin=*]
\item[4$^o$)] If
\[
c:a\longrightarrow b
\]
is an arbitrary path in $X$, we would like to deduce from this an isomorphism in Hot (of ``transport by $c$'') % [sic] isommorphisme -> isomorphisme, doubled m
\[
\overline{X'}_a\xrightarrow[\sim]{\bar{c}}\overline{X'}_b,
\]
with of course transitivity for composition,

\msp{page 141}

which will also give
\[
\overline{c^{\mathrm{opp}}}=\overline{c}^{-1}
\]
(assuming that the trivial paths $a\xrightarrow{\mathrm{id}}a$ give the identity $\overline{X'}_a\xrightarrow[\sim]{\mathrm{id}}\overline{X'}_a$!) (\footnotemark[32]\textsuperscript{,}\footnotemark[33]). Of course, $\overline{c}$ must depend only on the \textit{homotopy class} of $\overline{c}$ (in % [sic] de $\overline{c}$ (dans -> de $c$ (dans, the homotopy class is apparently that of the path c
\footnotetext[32]{Above all, not to confine oneself to the case where $c$ would be one of the $t_n$!}
\footnotetext[33]{describe the effect on the $\pi_0$ all the same!}
\printedpage{88}
% head: A. Grothendieck : Dérivateurs, Ch. VII, Édition des Universités de Montpellier II et Paris VII
% foot: Version 22 décembre 2001
$\uCh_\tau(X;a,b)$, or even in $C=\uCh(X;a,b)$, category of all paths in $X$, from $a$ to $b$, of unspecified type). In other words, we must have a functor
\[
\begin{array}{rcl}
\Pi(X) & \longrightarrow & \mathrm{Hot}\\[4pt]
a & \longmapsto & \hot(\overline{X'}_a)
\end{array}
\]
from the fundamental groupoid of $X$ to Hot. To tell the truth, this has nothing to do with the particular situation studied here - it must work whenever we have a HOT-fibration, here $\overline{X'}\longrightarrow X$, but also $\uCh_\infty(X)\longrightarrow X\times X$ (\footnote[34]{Here too, one should describe the effect of the transport $\overline{c}$ on the $\pi_0$ of the fibers.}). And for a $\mathbf{D}$-fibration $X'\longrightarrow X$ ($\mathbf{D}$ an arbitrary derivator), and $a\in\mathbf{D}(e)$, we must likewise have two functors
\[
\Pi(X)\rightrightarrows\mathbf{D}(e)
\]
corresponding to
\[
\left\{\begin{array}{lcl}
x & \longmapsto & H^\bullet_{\mathbf{D}}(X'_x,\bullet)\\[4pt]
x & \longmapsto & H_\bullet^{\mathbf{D}}(X'_x,\bullet).
\end{array}\right.
\]

\msp{page 142}

The case we are concerned with here would be the one where $\mathbf{D}=\mathrm{HOT}$, where we take $H^{\mathbf{D}}_\bullet$, and $a=$ final object of $\mathbf{D}(e)=\mathrm{Hot}$.)

\item[5$^o$)] If 1$^o$, 2$^o$, 3$^o$ are verified, we are morally sure that the construction of the $\uCh_\infty(X)$ is indeed the exact equivalent of the formalism of \textsc{Cartan-Serre}. To be completely sure of it, i.e. to link this theory explicitly to the classical constructions, one would have to prove a \textit{comparison statement} with semi-simplicial theory, which would essentially say this: If we denote, for $X$ in Cat, by $\tilde{X}$ the associated object in $\mathbf{\Delta}^\wedge$ (the ``nerve'' of $X$), then the arrow
\[
\uCh_\infty(X)\tilde{}\longrightarrow\underbrace{(X\times X)\tilde{}}_{\simeq\,\tilde{X}\times\tilde{X}}
\]
in $\mathbf{\Delta}^\wedge/(X\times X)\tilde{}$, is $W$-equivalent over $(X\times X)\tilde{}$ to the arrow
\[
\uCh(\tilde{X})\longrightarrow\tilde{X}\times\tilde{X}
\]
of the semi-simplicial theory of paths (assuming that there is such a classical arrow defined for \textit{every} object $X$ of $\mathbf{\Delta}^\wedge$). Or that we have at least canonical isomorphisms in Hot
\[
\hot_{\mathrm{Cat}}(\uCh_\infty(X;a,b))\ \simeq\ \hot_{\mathbf{\Delta}^\wedge}(\uCh(\tilde{X};\tilde{a},\tilde{b})).
\]
\end{enumerate}